% ==================================================
% CHAPTER 9: CONCLUSION
% ==================================================

\chapter{Conclusion}
\label{ch:conclusion}

\section{Summary}

MemoryLane turns the pages a person reads into a memory they can search by
describing what they remember. Over the semester it grew from a local,
single-user prototype into a deployed, multi-user product: a Chrome extension
that saves pages and videos, with automatic capture that keeps only what was
read and looks like content; a cleaning, chunking and embedding pipeline; a
hybrid search that fuses BM25 and vector retrieval by rank, understands sites
and times in queries, filters weak matches and shows readable snippets; and a
web app with a library, undoable deletes, summaries and voice search, all
behind Cognito sign-in with isolation enforced in the data layer, and running
at \texttt{memorylane.site}.

\section{Achievements}

\begin{enumerate}
  \item \textbf{Retrieval quality}: on SciFact the live search puts the right
  abstract first for 71\% of claims and in the top ten for 93\%, with MRR@10
  0.786, ahead of the prototype on every metric.
  \item \textbf{Measurement that changed the product}: the expanded
  evaluation found a ranking defect that raised Hit@1 from 0.44 to 0.71 once
  fixed, a gain confirmed by non-overlapping confidence intervals.
  \item \textbf{Speed}: about 76\,ms per search at the median on modest
  hardware, with no GPU.
  \item \textbf{Security and privacy by construction}: per-user isolation in
  the Solr client, PKCE sign-in, verified-email account linking, privacy-safe
  logs, and automatic capture that refuses private pages twice.
  \item \textbf{Engineering discipline}: 551 automated tests, regression
  tests for every post-launch bug, eleven decision records, and a 48-page
  documentation site with its own meaning-based search.
\end{enumerate}

\section{Limitations}

\begin{itemize}
  \item \textbf{Benchmark scope}: SciFact is a scientific benchmark, and the
  personal-library evaluation used small query sets; neither fully
  represents everyday browsing.
  \item \textbf{Single server}: one instance is a single point of failure
  and caps throughput; the embedding model runs on the CPU in one process.
  \item \textbf{Capture coverage}: pages that render little text, videos
  without captions and content behind logins are saved poorly or not at all.
  \item \textbf{Server-side storage}: page text is stored on the server, and
  summaries send a page to an external model when asked; users must trust
  the operator.
  \item \textbf{Chrome only}: the extension targets Chromium browsers and is
  installed unpacked rather than from a store.
\end{itemize}

\section{Future Work}

\begin{enumerate}
  \item \textbf{Larger and broader evaluation}: more BEIR datasets and a
  larger personal query set, run automatically after every retrieval change.
  \item \textbf{Learned reranking}: enable the existing cross-encoder stage
  where memory allows, and measure it.
  \item \textbf{Summary fallback}: a second, cheaper model or an extractive
  summary when Gemini is unavailable.
  \item \textbf{Import}: bring in existing bookmarks and browser history.
  \item \textbf{Account features}: password change and signing out of every
  device from an account page.
  \item \textbf{Local-first mode}: on-device embedding and storage for users
  who do not want page text on a server.
  \item \textbf{Store distribution}: publish the extension in the Chrome Web
  Store, and port it to Firefox.
\end{enumerate}

\section{Final Remarks}

The prototype showed that content-based search over one's own reading was
possible. The final product shows that it can be made dependable, private
and fast enough to use every day, and that careful measurement matters as
much as new features: the largest single improvement of the semester came
from evaluating the system honestly and fixing what the numbers revealed.
